\session{Lecture 1}{The cosmic microwave background}{I. Expansion, recombination, and sound waves}

\begin{frame}{Our route through the first two hours}
\begin{enumerate}
\item An expanding Universe: distances, energy, and temperature.
\item An opaque plasma becomes transparent.
\item Gravity and photon pressure create sound waves.
\item Recombination captures their phases.
\end{enumerate}
\key{The destination: understand where the acoustic pattern comes from.}
\medskip A 10-minute break near the middle. Questions and short notebook demonstrations throughout.
\end{frame}

\begin{frame}{From an unexpected signal to precision cosmology}
\begin{description}
\item[Late 1940s] Alpher and Herman predict a relic radiation bath from a hot early Universe, in the context of Gamow's hot-Big-Bang work.
\medskip
\item[1965] Penzias and Wilson publish their excess microwave antenna temperature. Dicke and collaborators identify the cosmological interpretation.
\medskip
\item[1990 / 1992] COBE establishes the nearly perfect blackbody spectrum (FIRAS), then detects primordial temperature anisotropies (DMR).
\medskip
\item[2000s / 2010s] WMAP and Planck map the sky in increasing detail: the tiny fluctuations become a precision test of cosmology.
\end{description}
\key{Two discoveries: a nearly uniform thermal background, and tiny variations across the sky.}
\credit{Historical background: \href{https://www.esa.int/Science_Exploration/Space_Science/Planck/Planck_and_the_cosmic_microwave_background}{ESA}; \href{https://lhea.gsfc.nasa.gov/archive/mather/cobe/history.html}{NASA/COBE history}; \href{https://lambda.gsfc.nasa.gov/education/lambda_graphics/cmb_discovery.html}{NASA/LAMBDA}.}
\end{frame}

\begin{frame}{First look at the sky: Planck's CMB temperature map}
\begin{columns}[c]
\column{.65\textwidth}
\centering\includegraphics[width=\linewidth,height=4.6cm,keepaspectratio]{planck-2018-cmb-map.jpg}
\column{.32\textwidth}
\small
\begin{itemize}
\item The whole sky, flattened into an oval (Mollweide projection).
\item False colour: orange/red is warmer; blue is cooler than the mean.
\item Tiny differences: typically $\sim100\,\mu$K on a $2.725$ K background.
\end{itemize}
\end{columns}
\key{What sizes of hot and cold patches are present? The angular power spectrum makes this quantitative.}
\credit{\copyright\ ESA/Planck Collaboration, \href{https://www.esa.int/ESA_Multimedia/Images/2018/07/Planck_s_view_of_the_cosmic_microwave_background2}{Planck Legacy Release 2018}. Observed, foreground-cleaned CMB reconstruction, not a simulation or a raw detector image. The mean and motion-induced dipole are removed.}
\end{frame}

\begin{frame}{This is the pattern we want to explain}
\plot{planck-notebook-1.png}{How much temperature variation is present at each angular scale? Larger multipoles correspond to smaller angles. Real Planck 2018 TT bandpowers and published best-fit template, from the Lecture 2 notebook; formal definitions come in Lecture 2.}
\end{frame}

\begin{frame}{A nearly uniform blackbody sky}
\[
B_\nu(T)=\frac{2h\nu^3}{c^2}\frac{1}{e^{h\nu/(k_BT)}-1},
\qquad T_0\simeq2.725\ \mathrm K.
\]
\[
T(\bn)=T_0[1+\Theta(\bn)],\qquad \Theta=\frac{\delta T}{T_0}.
\]
After removing the motion-induced dipole, typical temperature variations are of order $10^{-5}$.
\key{The mean temperature records thermal history. The small variations record gravity and the contents of the Universe.}
\credit{Convention: $c=\hbar=k_B=1$ unless units are explicitly restored.}
\end{frame}

\begin{frame}{Expansion changes physical distances}
\begin{columns}[c]
\column{.49\textwidth}
\[ds^2=-dt^2+a^2(t)d\bx^2\]
\[d_{\rm phys}=a(t)|\Delta\bx|\]
Comoving coordinates $\bx$ label the same freely moving observers. We choose $a(t_0)=1$.
\column{.48\textwidth}
\begin{tikzpicture}[scale=.8]
% Draw lines and dots from the same integer lattice in each local frame.
% TikZ's grid operation otherwise anchors its steps to the global origin.
\begin{scope}[x=.6cm,y=.6cm]
\foreach \i in {0,...,3}{
  \draw[teal!35] (\i,0)--(\i,3) (0,\i)--(3,\i);
  \foreach \j in {0,...,3}\fill[teal] (\i,\j) circle (1.8pt);
}
\end{scope}
\draw[-{Stealth},thick,orange] (2,.9)--(2.7,.9);
\begin{scope}[xshift=3cm,x=.9cm,y=.9cm]
\foreach \i in {0,...,3}{
  \draw[teal!35] (\i,0)--(\i,3) (0,\i)--(3,\i);
  \foreach \j in {0,...,3}\fill[teal] (\i,\j) circle (1.8pt);
}
\end{scope}
\node at (.9,-.5) {$a_1$};\node at (4.35,-.5) {$a_2>a_1$};
\end{tikzpicture}
\end{columns}
\ask{What stays fixed in this picture? What grows?}
\end{frame}

\begin{frame}{What controls the expansion?}
\[
H\equiv\frac{1}{a}\frac{da}{dt},\qquad
H^2=\frac{8\pi G}{3}\rho
\]
\[
\frac{1}{a}\frac{d^2a}{dt^2}=-\frac{4\pi G}{3}(\rho+3P).
\]
\begin{itemize}
\item Density determines the expansion rate.
\item Pressure also gravitates.
\item $H^{-1}$ gives the instantaneous expansion timescale; the age follows from the full expansion history.
\end{itemize}
\credit{Spatially flat FLRW background; $t$ is cosmic time.}
\end{frame}

\begin{frame}{Local conservation does not mean constant total energy}
\[
\nabla_\mu T^{\mu\nu}=0
\quad\Longrightarrow\quad
\frac{d\rho}{dt}+3H(\rho+P)=0.
\]
For a fixed comoving region, $V\propto a^3$ and $E=\rho V$:
\[dE=-P\,dV.\]
\key{A general expanding FLRW spacetime lacks the time-translation symmetry that would define a globally conserved energy of the usual kind.}
For a separately conserved fluid with $P=w\rho$,
\[\rho\propto a^{-3(1+w)},\qquad E\propto a^{-3w}.\]
\end{frame}

\begin{frame}{Energy inside the same comoving region}
\plot{local-energy-plot.pdf}{Each component is normalized to its own present energy. The curves compare how the components evolve; their relative heights do not show their abundances.}
\end{frame}

\begin{frame}{Different components dominate at different times}
\plot{plot-expansion.pdf}{Lecture 1 notebook. Matter--radiation equality is where $\Omega_r a^{-4}=\Omega_m a^{-3}$.}
\end{frame}

\begin{frame}{Expansion cools the photon bath}
\[
\lambda\propto a,\qquad E_\gamma\propto a^{-1},\qquad
1+z=\frac{1}{a}.
\]
Uniform redshifting preserves a Planck spectrum:
\[\boxed{T(z)=T_0(1+z).}\]
At $z\simeq1100$, $T\simeq3000$ K, or $k_BT\simeq0.26$ eV.
\ask{Why does radiation density fall as $a^{-4}$ rather than $a^{-3}$?}
\end{frame}

\begin{frame}{Define the distance before following the photon}
\begin{columns}[c]
\column{.57\textwidth}
\[d\tau=\frac{dt}{a},\quad \chi=|\bx|\]
For incoming radial light: $d\chi/d\tau=-1$.
\[\chi_* =\tau_0-\tau_* =\int_0^{z_*}\frac{dz}{H(z)}.\]
$L_{\rm com}$: transverse comoving size of a feature on the last-scattering shell.
\[\theta\simeq\frac{L_{\rm com}}{\chi_*}\qquad\text{(small angle)}\]
\column{.40\textwidth}
\begin{tikzpicture}[scale=.85]
\draw[very thick,teal] (0,0) circle (2);
\fill[navy] (0,0) circle(2pt);
\draw[orange,very thick] (20:2) arc(20:45:2);
\draw[navy] (0,0)--(20:2) (0,0)--(45:2);
\draw[orange] (20:.65) arc(20:45:.65);
\node at (-.15,-.35) {us};
\node at (1.1,.12) {$\chi_*$};
\node[orange] at (2.4,1.25) {$L_{\rm com}$};
\node at (0,-2.5) {last-scattering shell};
\end{tikzpicture}
\end{columns}
\key{Later, our physical ruler will be the sound horizon.}
\credit{From now on, dots denote $d/d\tau$ and $\Hs=\dot a/a=aH$. Distances above use $c=1$.}
\end{frame}

\begin{frame}{There is more than one useful horizon}
\[
H\equiv\frac1a\frac{da}{dt},\qquad
\boxed{\Hs\equiv\frac1a\frac{da}{d\tau}=aH.}
\]
$H$ uses cosmic time $t$; $\Hs$ uses conformal time $\tau$.
\medskip
\begin{description}
\item[Comoving Hubble length] $\Hs^{-1}$: the expansion timescale expressed as a comoving length.
\item[Comoving particle horizon] $\int_{t_i}^{t}dt'/a(t')$: how far light could have travelled.
\item[Sound horizon] $r_s(\tau)=\int_0^\tau c_s(\tau')d\tau'$: how far a pressure disturbance could have travelled.
\end{description}
\key{For acoustic oscillations, the relevant causal scale is the sound horizon.}
\end{frame}

\begin{frame}{Why was the early Universe opaque?}
\[\Gamma_T=n_e\sigma_T,\qquad \lambda_{\rm mfp}=(n_e\sigma_T)^{-1}.\]
\begin{center}
\begin{tikzpicture}[node distance=26mm, every node/.style={align=center}]
\node[draw=teal,rounded corners,minimum width=25mm,minimum height=12mm] (p) {photons};
\node[draw=teal,rounded corners,right=of p,minimum width=25mm,minimum height=12mm] (e) {electrons};
\node[draw=teal,rounded corners,right=of e,minimum width=25mm,minimum height=12mm] (b) {ions};
\draw[<->,thick] (p)--node[above=2mm,font=\small]{Thomson}(e);
\draw[<->,thick] (e)--node[above=2mm,font=\small]{Coulomb}(b);
\end{tikzpicture}
\end{center}
\[\Gamma\gg H:\ \text{frequent interactions};\qquad\Gamma\lesssim H:\ \text{decoupling}.\]
\ask{What happens to the opacity when free electrons disappear into atoms?}
\end{frame}

\begin{frame}{How does equilibrium divide hydrogen between ions and atoms?}
\[p+e^-\leftrightarrow H^0+\gamma,\qquad \mu_p+\mu_e=\mu_{H^0},\quad\mu_\gamma=0.\]
For a dilute non-relativistic species, including rest energy in $\mu_i$,
\[n_i=g_i\left(\frac{m_iT}{2\pi}\right)^{3/2}e^{(\mu_i-m_i)/T}.\]
\begin{itemize}
\item $n_e$ and $n_p$ count \emph{free} electrons and protons.
\item $n_{H^0}$ counts neutral hydrogen atoms, a separate species.
\item $g_i$ counts the internal states of species $i$.
\end{itemize}
\key{Form $n_en_p/n_{H^0}$: chemical equilibrium will eliminate the chemical potentials.}
\credit{Hydrogen only; natural units. The full derivation is in Section 5.2 of the notes.}
\end{frame}

\begin{frame}{Use equilibrium to relate the three species}
Insert the three species densities:
\[
\frac{n_en_p}{n_{H^0}}=
\frac{g_eg_p}{g_{H^0}}
\left(\frac{m_em_p}{m_{H^0}}\frac{T}{2\pi}\right)^{3/2}
e^{(\mu_e+\mu_p-\mu_{H^0})/T}
e^{-(m_e+m_p-m_{H^0})/T}.
\]
\begin{itemize}
\item Chemical equilibrium: $\mu_e+\mu_p-\mu_{H^0}=0$.
\item Binding energy: $m_e+m_p-m_{H^0}=E_I=13.6$ eV.
\item In the prefactor only, $m_{H^0}\simeq m_p$; ground-state spin counts give $g_eg_p/g_{H^0}=2\times2/4=1$.
\end{itemize}
\[
\boxed{\frac{n_en_p}{n_{H^0}}\simeq
\left(\frac{m_eT}{2\pi}\right)^{3/2}e^{-E_I/T}.}
\]
\credit{The tiny binding-energy mass difference must be retained in the exponential, where it supplies the temperature dependence.}
\end{frame}

\begin{frame}{Count every hydrogen nucleus}
\[
\boxed{n_H=\underbrace{n_p}_{\text{free protons}}+
\underbrace{n_{H^0}}_{\text{neutral atoms}}}
\]
Each neutral atom contains one proton: $n_H$ counts \emph{all} hydrogen nuclei.
\par Define the ionized fraction using this total; charge neutrality gives $n_e=n_p$:
\[
x_e=\frac{n_e}{n_H},\qquad n_e=n_p=x_en_H,\qquad
n_{H^0}=(1-x_e)n_H.
\]
Substitute into the ratio:
\[
\frac{n_en_p}{n_{H^0}}=
\frac{(x_en_H)^2}{(1-x_e)n_H}=n_H\frac{x_e^2}{1-x_e}.
\]
Divide the preceding Saha ratio by $n_H$ to obtain an equation for $x_e$.
\credit{Only in the fully ionized limit is $n_H\simeq n_p$. The ratio $n_e/n_p=1$ follows from charge neutrality and therefore does not measure ionization.}
\end{frame}

\begin{frame}{Why recombination occurs far below 13.6 eV}
Combining the equilibrium ratio with the ionized fraction gives the \emph{Saha equation}:
\[
\boxed{\frac{x_e^2}{1-x_e}=\frac1{n_H}
\left(\frac{m_eT}{2\pi}\right)^{3/2}e^{-E_I/T}.}
\]
\begin{itemize}
\item At $T=E_I$, the density-dependent prefactor is enormous: $x_e\simeq1$.
\item Roughly a billion photons per baryon keep the ionizing tail populated.
\item Half ionization requires $E_I/T\simeq42$, giving $T\simeq0.32$ eV.
\end{itemize}
\key{This equilibrium result determines the ionization fraction. Last scattering also depends on the interaction rate and expansion.}
\end{frame}

\begin{frame}{Neutral hydrogen forms rapidly as the plasma cools}
\plot{saha-calculation.pdf}{Hydrogen Saha equilibrium. The cosmological evolution proceeds from right to left on this redshift axis.}
\end{frame}

\begin{frame}{Last scattering is a probability distribution}
\[
\widetilde\tau(\tau)=\int_\tau^{\tau_0} a n_e\sigma_T\,d\tau',
\qquad \dot{\widetilde\tau}<0.
\]
\begin{itemize}
\item $e^{-\widetilde\tau}$ is the probability of reaching us without another scattering.
\item The visibility is $\boxed{g_\tau=-\dot{\widetilde\tau}e^{-\widetilde\tau}}$.
\item The finite visibility width corresponds to a last-scattering shell with finite thickness.
\end{itemize}
\credit{The tilde matters: $\widetilde\tau$ is optical depth; $\tau$ is conformal time.}
\end{frame}

\begin{frame}{Recombination and last scattering are related, but distinct}
\plot{visibility-calculation.pdf}{Saha-model demonstration. $g_z=g_\tau/H(z)$, so the peak depends on the integration variable. Precision calculations place last scattering near $z_*\simeq1100$; this equilibrium toy model peaks earlier.}
\end{frame}

\begin{frame}{Ten-minute break}
\vfill\Large We have an expanding, cooling Universe that becomes transparent.
\medskip\par What could make different patches of the last-scattering shell have different temperatures?
\vfill
\end{frame}

\begin{frame}{Add small perturbations to the smooth background}
\[\rho_i(\tau,\bx)=\bar\rho_i(\tau)[1+\delta_i(\tau,\bx)].\]
In conformal Newtonian gauge,
\[ds^2=a^2\big[-(1+2\Psi)d\tau^2+(1-2\Phi)d\bx^2\big].\]
\begin{itemize}
\item $\Psi$: gravitational potential; a potential well has $\Psi<0$.
\item $\Phi$: scalar perturbation of spatial curvature.
\item Keep terms linear in potentials, density contrasts, and velocities.
\end{itemize}
\end{frame}

\begin{frame}{Why do two scalar potentials remain?}
\begin{center}\Large
4 scalar metric functions\par\medskip
$-$ 2 scalar coordinate freedoms\par\medskip
$=$ 2 scalar metric potentials
\end{center}
\key{These are constrained scalar potentials. The two propagating gravitational-wave polarizations belong to the tensor sector.}
Einstein's equations relate the scalar potentials to matter perturbations.
\end{frame}

\begin{frame}{Anisotropic stress is directional momentum flux}
\[T^i{}_j=P\delta^i{}_j+\pi^i{}_j,\qquad \pi^i{}_i=0.\]
\begin{itemize}
\item A quadrupole of particle momenta gives stress beyond isotropic pressure.
\item Free-streaming neutrinos arrive from different regions without frequent collisions to isotropize them.
\item Tight coupling suppresses photon anisotropic stress; photons develop it as scattering weakens.
\end{itemize}
\key{Negligible total scalar anisotropic stress gives $\Phi=\Psi$. Cold matter has small random-velocity stresses.}
\end{frame}

\begin{frame}{Why switch from position $\bx$ to wave number $\bk$?}
\small
The perturbation contains many spatial scales. Because the equations are linear, we may evolve one sinusoidal component at a time:
\begin{columns}[c]
\column{.48\textwidth}
\[\delta(\bx)=\int\frac{d^3k}{(2\pi)^3}\delta(\bk)e^{i\bk\cdot\bx}.\]
\begin{itemize}
\item $\bk$ gives the direction; $k=|\bk|$ gives the scale.
\item The wavelength is $\lambda=2\pi/k$: large $k$ means a short wave.
\item Derivatives become algebraic: $\nabla\to i\bk$ and $\nabla^2\to-k^2$.
\end{itemize}
\column{.49\textwidth}\includegraphics[width=\linewidth]{spatial-wave.pdf}
\end{columns}
\key{The drawing is one mode at one instant: crests are compressions and troughs are rarefactions. We will solve for how its amplitude changes with time.}
\end{frame}

\begin{frame}{Plan for deriving the acoustic wave}
For one Fourier mode, we track
\[
\underbrace{\Theta_0=\delta_\gamma/4}_{\text{photon monopole: density}},\qquad
\underbrace{\theta_i=i\bk\!\cdot\!\bv_i\ (i=\gamma,b)}_{\text{velocity divergence}},\qquad
\underbrace{\Phi,\Psi}_{\text{gravity}}.
\]
For photons, $\theta_\gamma$ carries the same bulk-flow information as the temperature dipole $\Theta_1$.
\begin{columns}[T]
\column{.47\textwidth}
\begin{block}{Two evolution equations}
\textbf{Continuity:} motion changes density.\\[1mm]
\textbf{Euler:} pressure and gravity change motion.
\end{block}
\column{.49\textwidth}
\begin{block}{Scattering and baryons}
Scattering drives $\theta_\gamma$ toward $\theta_b$. Baryons then add inertia to the pressure-supported photon fluid.
\end{block}
\end{columns}
\key{We will combine continuity and Euler into one oscillator equation. Pressure restores, gravity drives, and baryons change the inertia.}
\end{frame}

\begin{frame}{Velocity moves the fluid; convergence compresses it}
\vspace{-1.5mm}
\begin{columns}[T]
\column{.47\textwidth}
\begin{block}{Uniform translation}
\centering
\begin{tikzpicture}[x=1cm,y=1cm]
  \draw[rounded corners,fill=cyan!6,draw=teal!65] (0,0) rectangle (4.3,1.55);
  \foreach \x/\y in {.55/.40,1.35/1.10,2.05/.55,2.85/1.12,3.65/.45}
    \fill[teal] (\x,\y) circle (2.2pt);
  \foreach \y in {.35,.78,1.20}
    \draw[-{Latex[length=2.5mm]},very thick,orange] (-.15,\y) -- (.65,\y);
  \foreach \y in {.35,.78,1.20}
    \draw[-{Latex[length=2.5mm]},very thick,orange] (3.65,\y) -- (4.45,\y);
\end{tikzpicture}

The same amount enters and leaves:
\[
\bv_\gamma\neq0,qquad \nabla\!\cdot\bv_\gamma=0.
\]
The patch moves, but its density is unchanged.
\end{block}

\column{.49\textwidth}
\begin{block}{Converging flow}
\centering
\begin{tikzpicture}[x=1cm,y=1cm]
  \draw[rounded corners,fill=cyan!12,draw=teal!65] (0,0) rectangle (4.3,1.55);
  \foreach \x/\y in {1.45/.35,1.75/.78,1.48/1.20,2.20/.48,2.48/.95,2.85/.38,2.82/1.22}
    \fill[teal] (\x,\y) circle (2.2pt);
  \foreach \y in {.35,.78,1.20}{
    \draw[-{Latex[length=2.5mm]},very thick,orange] (-.15,\y) -- (.75,\y);
    \draw[-{Latex[length=2.5mm]},very thick,orange] (4.45,\y) -- (3.55,\y);
  }
\end{tikzpicture}

More enters than leaves:
\[
\theta_\gamma\equiv\nabla\!\cdot\bv_\gamma<0.
\]
Photons accumulate, so the patch becomes denser and hotter.
\end{block}
\end{columns}
\key{Continuity counts net flow through a boundary: $\dot\Theta_0=-\theta_\gamma/3+\dot\Phi$. The velocity matters through its spatial variation.}
\end{frame}

\begin{frame}{Baryon loading: photon pressure moves two fluids}
Momentum density is proportional to $(\rho+P)\bv$.
\[
\boxed{R=\frac{\bar\rho_b}{4\bar\rho_\gamma/3}}
\quad\Longrightarrow\quad
\text{total inertia}=\frac{4\bar\rho_\gamma}{3}(1+R).
\]
Photons supply the pressure while photons and baryons both have to accelerate. For an adiabatic compression, $\delta\rho_b=R\delta\rho_\gamma$, hence
\[c_s^2=\frac{\delta P_\gamma}{\delta\rho_\gamma+\delta\rho_b}
=\boxed{\frac1{3(1+R)}}.\]
\key{$R$ is the extra baryon inertia measured in units of photon inertia. Larger $R$ lowers the sound speed.}
\end{frame}

\begin{frame}{As the Universe expands, baryons become harder to move}
\plot{sound-horizon-calculation.pdf}{The left panel shows $R=3\bar\rho_b/(4\bar\rho_\gamma)$; the right shows $c_s/c=[3(1+R)]^{-1/2}$. Cosmic time runs from high redshift on the right toward low redshift on the left.}
\end{frame}

\begin{frame}{Step 1: converging motion compresses the photons}
\[\boxed{\dot\Theta_0=-\frac13\theta_\gamma+\dot\Phi.}\]
\begin{itemize}
\item $\Theta_0=\delta_\gamma/4$ measures the photon-temperature monopole.
\item $\theta_\gamma=i\bk\cdot\bv_\gamma<0$ means that the photon flow converges, so $\Theta_0$ rises.
\item $\dot\Phi$ accounts for compression caused by changing spatial curvature.
\end{itemize}
The equivalent density equation is $\dot\delta_\gamma=-4\theta_\gamma/3+4\dot\Phi$.
\ask{With constant $\Phi$, what sign of $\theta_\gamma$ makes a patch hotter?}
\end{frame}

\begin{frame}{Step 2: pressure and gravity set photons in motion}
\[
\dot\theta_\gamma=
\underbrace{k^2\Theta_0}_{\text{pressure}}
+\underbrace{k^2\Psi}_{\text{gravity}}
+\underbrace{\Gamma_c(\theta_b-\theta_\gamma)}_{\text{scattering}},
\]
\[\Gamma_c\equiv-\dot{\widetilde\tau}=a n_e\sigma_T>0.\]
\begin{itemize}
\item A temperature crest has higher photon pressure, which pushes material away from it.
\item Gravity pulls the fluid into potential wells.
\item Scattering acts whenever the photon and baryon flows differ.
\end{itemize}
\credit{There is no photon Hubble drag: $\Hs(1-3w_\gamma)\theta_\gamma=0$ because $w_\gamma=1/3$. The signs look different from vector forces because $\theta$ is a divergence. Leading tight coupling neglects the photon quadrupole; $\Gamma_c=a\Gamma_T$.}
\end{frame}

\begin{frame}{Step 3: scattering transfers momentum between the fluids}
\vspace{-2mm}
\[
\dot\theta_b+\underbrace{\Hs\theta_b}_{\text{expansion drag}}
=k^2\Psi+\frac{\Gamma_c}{R}(\theta_\gamma-\theta_b).
\]
\credit{Baryon pressure $c_b^2k^2\delta_b$ is negligible because $c_b^2\sim T_b/m_b\ll1$; photon pressure is transmitted by scattering.}
The photon--baryon \emph{slip} measures their relative flow:
\[\boxed{\Delta\theta\equiv\theta_\gamma-\theta_b}.\]
The collision terms exchange momentum internally:
\[\left.\dot\theta_\gamma\right|_{\rm coll}
+R\left.\dot\theta_b\right|_{\rm coll}=0.\]
\key{Scattering reduces the slip. The same force density gives a baryon acceleration proportional to $1/R$ because $R$ measures baryon inertia.}
\end{frame}

\begin{frame}{Tight coupling means that the slip rapidly decays}
Keeping only the collision terms gives
\[
\frac{d\Delta\theta}{d\tau}=-\Gamma_c(1+R^{-1})\Delta\theta.
\]
\begin{center}
  \includegraphics[width=.66\linewidth,height=.37\textheight,keepaspectratio]{collision-demonstration.pdf}
\end{center}
\credit{Here $u=\Gamma_c(\tau-\tau_i)$ counts elapsed scattering times in a fixed-$R$, fixed-rate collision-only example. The two flows approach their inertia-weighted common value.}
\key{After this short transient, $\theta_\gamma\simeq\theta_b\equiv\theta$. We can then describe one tightly coupled fluid.}
\end{frame}

\begin{frame}{Step 4: one momentum equation for the coupled plasma}
\small
Start from the two equations, weighting the baryon equation by its inertia:
\[
\begin{aligned}
\text{photon Euler:}\quad
\dot\theta_\gamma
 &=k^2\Theta_0+k^2\Psi
   +\Gamma_c(\theta_b-\theta_\gamma),\\
R\times\text{baryon Euler:}\quad
R\dot\theta_b+\Hs R\theta_b
 &=Rk^2\Psi
   +\Gamma_c(\theta_\gamma-\theta_b).
\end{aligned}
\]
Add them: the collision terms cancel. With
$\theta_\gamma\simeq\theta_b\equiv\theta$,
\[(1+R)\dot\theta+\Hs R\theta=k^2\Theta_0+(1+R)k^2\Psi.\]
Using $\dot R=\Hs R$ and dividing by $1+R$ gives
\[\boxed{\dot\theta+\frac{\dot R}{1+R}\theta
=\frac{k^2}{1+R}\Theta_0+k^2\Psi.}\]
\key{Photon pressure is divided by the total inertia $1+R$, while gravity accelerates both components.}
\credit{The slip is small, but $\Gamma_c(\theta_\gamma-\theta_b)$ remains the finite internal force that transfers photon momentum to baryons.}
\end{frame}

\begin{frame}{Step 5: continuity and Euler give an acoustic oscillator}
Differentiate continuity and use the combined momentum equation to eliminate $\theta$:
\[
\boxed{\ddot\Theta_0+
\underbrace{\frac{\dot R}{1+R}\dot\Theta_0}_{\text{changing inertia}}+
\underbrace{k^2c_s^2\Theta_0}_{\text{pressure restoring force}}
=\underbrace{S_{\rm grav}}_{\text{gravitational driving}}}
\]
\[S_{\rm grav}=\ddot\Phi+\frac{\dot R}{1+R}\dot\Phi-\frac{k^2}{3}\Psi.\]
\begin{itemize}
\item The natural frequency is approximately $k c_s$: short waves oscillate faster.
\item Baryon inertia lowers $c_s$ and changes the amplitude.
\item Constant or evolving potentials displace or drive the oscillation.
\end{itemize}
\end{frame}

\begin{frame}{Why introduce an effective temperature?}
The emitting plasma has intrinsic perturbation $\Theta_0$.
\medskip\par A photon leaving a static potential acquires fractional shift $\Psi_*-\Psi_{\rm obs}$.
\medskip\par The observer term is common to all directions. The emission contribution is
\[\boxed{X=\Theta_0+\Psi.}\]
\key{$X$ is the temperature perturbation at emission after including the gravitational redshift from the emitting potential.}
Doppler shifts and changing potentials along the path will enter Lecture 2.
\end{frame}

\begin{frame}{Make the driving and equilibrium explicit}
\small
\[
\begin{aligned}
\ddot X+\frac{\dot R}{1+R}\dot X+k^2c_s^2X
={}&\ddot\Phi+\ddot\Psi+\frac{\dot R}{1+R}(\dot\Phi+\dot\Psi)\\
&{}-k^2c_s^2R\Psi.
\end{aligned}
\]
\normalsize Freeze $R,\Phi,\Psi$ for a simple solvable model:
\[\boxed{\ddot X+k^2c_s^2(X+R\Psi)=0.}\]
\key{The force vanishes at $X=-R\Psi$, so baryons shift the centre of the oscillation away from zero.}
\credit{Freezing coefficients is a toy approximation. Expansion-related amplitude evolution is distinct from dissipative diffusion damping.}
\end{frame}

\begin{frame}{Derive the equilibrium from force balance}
No velocity or acceleration means pressure balances gravity:
\[0=\frac{k^2}{1+R}\Theta_{0,\rm eq}+k^2\Psi.\]
\[\Theta_{0,\rm eq}=-(1+R)\Psi,\qquad
\boxed{X_{\rm eq}=-R\Psi.}\]
\begin{itemize}
\item Photons must support their own inertia \emph{and} the baryons.
\item For $R=0$, intrinsic temperature and gravitational redshift cancel at equilibrium.
\item For $\Psi<0$ and $R>0$, the effective-temperature offset is positive.
\end{itemize}
\end{frame}

\begin{frame}{How to read the baryon-loading toy model}
\small
The vertical axis divides by $|\Psi|$ to remove the arbitrary amplitude of this linear mode.
\[
s=k(\tau-\tau_i),\qquad
\varphi=k r_s=\frac{s}{\sqrt{3(1+R)}}.
\]
\begin{center}
\includegraphics[width=.78\linewidth,height=.36\textheight,keepaspectratio]{oscillator-demonstration.pdf}
\end{center}
\begin{itemize}
\item \textbf{Left, same $s$:} larger $R$ oscillates more slowly.
\item \textbf{Right, same $\varphi$:} aligned cycles expose the shifted equilibrium and unequal extrema.
\end{itemize}
\credit{$s$ measures elapsed time in wavelength units; $\varphi=kr_s$ is the acoustic phase derived from the sound speed. Lines solve one Fourier mode; circles show the analytic solution.}
\end{frame}

\begin{frame}{The sound horizon keeps track of acoustic phase}
\[r_s(\tau)=\int_0^\tau c_s\,d\tau',\qquad
X+R\Psi=C_1\cos(kr_s)+C_2\sin(kr_s).\]
\begin{itemize}
\item A displaced fluid with negligible initial velocity starts in the cosine phase.
\item Adiabatic growing modes have coherent temporal phases.
\item Recombination takes a snapshot at $r_s(\tau_*)$.
\end{itemize}
\key{Extrema satisfy $k_n r_s(\tau_*)\simeq n\pi$: a harmonic series.}
\credit{The expression is exact for frozen coefficients; slowly evolving coefficients also change the amplitude and equilibrium.}
\end{frame}

\begin{frame}{Temperature and velocity are out of phase}
\plot{acoustic-phases-plot.pdf}{The horizontal coordinate is acoustic phase $k r_s$. The curves show normalized phase relations; their vertical scales do not represent equal observable amplitudes.}
\end{frame}

\begin{frame}{Projection and diffusion set the final angular pattern}
\[\theta_s=\frac{r_s(\tau_*)}{\chi_*},\qquad \ell_A\simeq\frac{\pi}{\theta_s}.\]
\begin{itemize}
\item Sound supplies a ruler; the distance turns it into an angle.
\item Baryon loading modulates compression versus rarefaction peaks.
\item Potentials decaying near equality boost early-entering modes: radiation driving.
\item Photon diffusion damps short waves: amplitude $\sim e^{-k^2/k_D^2}$.
\item Finite last-scattering width averages over nearby phases.
\end{itemize}
\credit{$\ell_A$ gives the characteristic spacing. Potential evolution and projection shift the individual peaks, including the first.}
\end{frame}

\begin{frame}{Explore the physics in the notebook}
\Large\colabone\normalsize
\medskip
\begin{enumerate}
\item Compare 13.6 eV with the Saha half-ionization temperature.
\item Predict the final common velocity when $R$ changes.
\item Increase baryon loading: what happens to frequency and equilibrium?
\end{enumerate}
\key{Make a prediction before rerunning each example.}
\credit{All plots are included in these slides, so the session also works without an internet connection.}
\end{frame}

\begin{frame}{The physical story in five steps}
\begin{enumerate}
\item Expansion cools the photon bath.
\item Neutral-atom formation reduces opacity.
\item Gravity and photon pressure create acoustic oscillations.
\item Baryons shift equilibrium; diffusion damps short scales.
\item Last scattering projects a snapshot onto our sky.
\end{enumerate}
\key{Next: derive the angular power spectrum and measure it with real data.}
\end{frame}
